\documentclass[11pt]{article}
\usepackage[T1]{fontenc}
\usepackage{lmodern,microtype,amsmath,amssymb,amsthm,mathtools,mathrsfs,tikz-cd}
\usepackage[margin=1in]{geometry}
\usepackage{xurl,hyperref}
\hypersetup{colorlinks=true,linkcolor=blue,citecolor=blue,urlcolor=blue,
pdftitle={Inseparable domains and computational extensions},pdfauthor={Florian Lengyel}}
\newtheorem{theorem}{Theorem}[section]
\newtheorem{proposition}[theorem]{Proposition}
\newtheorem{lemma}[theorem]{Lemma}
\newtheorem{corollary}[theorem]{Corollary}
\newcommand{\R}{\mathscr R}
\newcommand{\D}{\mathscr D}
\newcommand{\Pair}{\operatorname{Pair}}
\newcommand{\Inv}{\operatorname{Inv}}
\newcommand{\Tot}{\operatorname{Tot}}
\newcommand{\ran}{\operatorname{ran}}
\title{Inseparable domains and computational extensions}
\author{Florian Lengyel}
\date{29 October 2026}
\begin{document}
\maketitle
\begin{abstract}
An inseparable pair of computational domains determines a consistent
presentation of signed literals on a binary coproduct. Its consistent
computational extensions are exactly the disjoint enlargements of the
original pair. Consequently, inseparability is equivalent to essential
incompleteness for this presentation. The construction is left adjoint
to the functor taking an involutive presentation to its positive and
negative domains. For classical theories, the existing atom constructor
transports the same obstruction to every consistent extension, whose
theorem and refutation domains are themselves inseparable. We also
separate reflection of inseparability by restriction functors from its
preservation; a projection of recursion categories disproves unrestricted
preservation, even when it preserves a universal evaluator.
\end{abstract}

\section{Domains, separators, and reductions}

Let $\R$ be a restriction category with restriction zeros.
Composition is $gf=g\circ f$. For $f:X\to Y$, its restriction
idempotent is $\bar f:X\to X$; $f$ is total when $\bar f=1_X$.
Write $\D(X)$ for the poset of restriction idempotents on $X$,
ordered by $e\le d$ when $de=e$. Its meet is composition, its bottom
is $0_X$, and its top is $1_X$. For a total arrow $f:X\to Y$, put
\[
 f^*:\D(Y)\longrightarrow\D(X),\qquad f^*e=\overline{ef}.
\]
The restriction identities give preservation of meets and both bounds,
as well as $(gf)^*=f^*g^*$ and $1_X^*=1_{\D(X)}$.
Thus domains and inverse images define
$\D:\Tot(\R)^{\mathrm{op}}\to\mathbf{Pos}$, where
$\mathbf{Pos}$ has posets as objects and monotone maps as arrows.
Domain idempotents remain arrows of $\R$; when split,
their splitting objects belong to its total category.

Whenever the domain poset has binary joins, a pair $(c,d)$ is
\emph{complementary} if $cd=0$ and $c\vee d=1$.
It separates $(a,b)$ if $a\le c$ and $b\le d$.
A disjoint pair $(a,b)$ is \emph{inseparable} if it has no such
separator.

Define the category $\Pair(\R)$ to have objects $(X,a,b)$ with
$a,b\in\D(X)$ and $ab=0_X$. A morphism
\[
 f:(X,a,b)\longrightarrow(Y,c,d)
\]
is a total arrow $f:X\to Y$ satisfying $a\le f^*c$ and $b\le f^*d$.
Identities and composition are inherited from $\R$:
if $c\le g^*e$, then $a\le f^*c\le(gf)^*e$, and similarly for
the second coordinate. These arrows are simultaneous reductions of
the two domains.

\begin{proposition}\label{prop:transport}
Suppose all $\D(X)$ have binary joins and total inverse images preserve
them. If $f:(X,a,b)\to(Y,c,d)$ is a morphism of $\Pair(\R)$ and
$(a,b)$ is inseparable, then $(c,d)$ is inseparable.
\end{proposition}
\begin{proof}
A complementary separator $(r,s)$ of $(c,d)$ pulls back to the
complementary pair $(f^*r,f^*s)$. The inequalities
$a\le f^*c\le f^*r$ and $b\le f^*d\le f^*s$ contradict
inseparability of $(a,b)$.
\end{proof}

\section{The free involutive presentation}

Assume $\R$ has binary restriction coproducts: these are coproducts
$X+Y$ whose injections $i_0:X\to X+Y$ and $i_1:Y\to X+Y$ are
total. The copair of $f:X\to Z$ and $g:Y\to Z$ is written $[f,g]$.
For arrows $f:X\to X'$ and $g:Y\to Y'$, their coproduct is
$f+g=[i'_0f,i'_1g]:X+Y\to X'+Y'$, where $i'_0,i'_1$
are the target injections. In particular, for domains $a$ on $X$
and $b$ on $Y$, write $a+b=[i_0a,i_1b]$.

\begin{lemma}\label{lem:decompose}
There is an order isomorphism
\[
 \D(X+Y)\cong\D(X)\times\D(Y),\qquad
 e\longmapsto(i_0^*e,i_1^*e),\qquad (a,b)\longmapsto a+b.
\]
It preserves all meets and joins that exist.
\end{lemma}
\begin{proof}
For a restriction idempotent $e$ on $X+Y$, the restriction identities
give $ei_j=i_j\overline{ei_j}$. Hence
$e=(i_0^*e)+(i_1^*e)$ by the coproduct property.
Conversely, $a+b$ is a restriction idempotent: the identities for its
restriction, precomposed with either injection, give
$\overline{a+b}\,i_0=i_0a$ and
$\overline{a+b}\,i_1=i_1b$. Its two inverse images are $a,b$.
Finally, $(a+b)(c+d)=ac+bd$, so the two maps preserve and reflect
order. The assertion about meets and joins follows from their order
universal properties.
\end{proof}

An \emph{involutive presentation} is a triple $(S,\nu,t)$ consisting
of a total involution $\nu:S\to S$, a domain $t\in\D(S)$, and
the consistency equation $t\,\nu^*t=0_S$.
Regard $S$ as the object of assertions, $\nu$ as their negation
constructor, and $t$ as the theorem domain. In particular, $\nu^*$ is monotone
inverse image, not complementation of domains.
Let $\Inv(\R)$ have these presentations as objects. A morphism
\[
 h:(S,\nu,t)\longrightarrow(T,\mu,u)
\]
is a total arrow $h:S\to T$ with $h\nu=\mu h$ and $t\le h^*u$.
Identity and composition preserve both conditions.
There is a covariant functor
\[
 J:\Inv(\R)\longrightarrow\Pair(\R),\qquad
 J(S,\nu,t)=(S,t,\nu^*t),\quad J(h)=h.
\]
Indeed, equivariance gives $\nu^*h^*=h^*\mu^*$, which verifies
the second domain inequality.

\begin{theorem}\label{thm:free}
The functor $J$ has a left adjoint $\Sigma$ given by
\[
 \Sigma(X,a,b)=(X+X,\sigma_X,a+b),\qquad
 \sigma_X=[i_1,i_0],\qquad \Sigma(f)=f+f.
\]
Neither domain joins nor ranges are needed for this adjunction.
\end{theorem}
\begin{proof}
The map $\sigma_X$ is a total involution and
$\sigma_X^*(a+b)=b+a$. Thus
$(a+b)\sigma_X^*(a+b)=ab+ba=0_{X+X}$.
For a morphism $f$ of pairs, $(f+f)^*(c+d)=f^*c+f^*d$;
this proves the required theorem-domain inequality. The coproduct
identities prove equivariance and the functor laws.

For $(S,\nu,t)\in\Inv(\R)$, the bijection defining the adjunction is
\[
 \Inv(\R)(\Sigma(X,a,b),(S,\nu,t))
 \cong\Pair(\R)((X,a,b),J(S,\nu,t)).
\]
It sends $h$ to $hi_0$. Its inverse sends $f$ to $[f,\nu f]$.
Every equivariant $h$ satisfies $hi_1=\nu hi_0$, so these assignments
are inverse. By Lemma~\ref{lem:decompose}, the inequality
$a+b\le[f,\nu f]^*t$ is equivalent to
$a\le f^*t$ and $b\le f^*\nu^*t$.
Both assignments commute with precomposition and postcomposition;
for postcomposition this uses the equivariance of the composed arrow.
The unit is $i_0:X\to X+X$, and the counit is
$[1_S,\nu]:S+S\to S$.
\end{proof}

For a fixed involutive presentation $(S,\nu,t_0)$, a
\emph{consistent computational extension} is a selected domain
$t\in\D(S)$ with $t_0\le t$ and $t\,\nu^*t=0_S$.
It is \emph{complete} if $t\vee\nu^*t=1_S$.
Essential incompleteness below refers to this specified class of
domain extensions.

\begin{theorem}\label{thm:signed}
Let $\R$ have restriction zeros and binary restriction coproducts,
and suppose $\D(X)$ has binary joins. For disjoint $a,b\in\D(X)$,
the following are equivalent:
\begin{enumerate}
\item $(a,b)$ is inseparable;
\item every consistent computational extension of
      $(X+X,\sigma_X,a+b)$ is incomplete.
\end{enumerate}
The base presentation is consistent. Its extension poset is isomorphic
to the coordinatewise ordered poset of pairs $(c,d)$ with
$a\le c$, $b\le d$, and $cd=0_X$.
\end{theorem}
\begin{proof}
Lemma~\ref{lem:decompose} writes each extension uniquely as $t=c+d$.
The three conditions are
\[
 \begin{aligned}
 a+b\le t&\ \Longleftrightarrow\ a\le c\text{ and }b\le d,\\
 t\,\sigma_X^*t=0&\ \Longleftrightarrow\ cd=0,\\
 t\vee\sigma_X^*t=1&\ \Longleftrightarrow\ c\vee d=1_X.
 \end{aligned}
\]
Thus a consistent complete extension is exactly a complementary
separator of $(a,b)$. Consistency of the base follows from $ab=0$.
\end{proof}

This construction applies within a recursion category as soon as its
inseparable pair and binary restriction coproducts are supplied.
It presents signed literals with involutive negation. Constructing
implication, the classical logical axioms, and modus ponens requires a
further deductive presentation. In particular, the theorem does not
supply a conservative extension to that calculus. Its domain calculation
uses no iteration, regularity, decidable equality, or proof algebra;
stability of arbitrary total inverse images is also unnecessary here,
because the coproduct decomposition computes the required joins.

\section{The fixed interpretation for classical extensions}

The same separation argument applies to the classical presentations of
\cite{Domains,General,Iteration}. Suppose $\R$ has restriction zeros,
binary domain joins preserved by total inverse images, an inseparable
pair $a,b\in\D(U)$, total arrows
\[
 \theta:U\to S,\qquad \nu:S\to S,
\]
and a domain $t_0\in\D(S)$ such that
\begin{equation}\label{eq:base}
 a\le\theta^*t_0,\qquad b\le\theta^*\nu^*t_0.
\end{equation}
Here $S$ is the syntax object, $\theta$ its atom constructor, and
$\nu$ its negation constructor; involutivity of $\nu$ is not required.

Let $\mathcal E(t_0)$ be the poset category of domains $t\ge t_0$
with $t\,\nu^*t=0_S$, with one arrow $t\to t'$ when $t\le t'$.
It may instead be restricted to the rule-closed computational theory
domains. There is a covariant functor
\[
 \mathcal P:\mathcal E(t_0)\longrightarrow\Pair(\R),\qquad
 t\longmapsto(S,t,\nu^*t),\qquad (t\le t')\longmapsto1_S.
\]
Monotonicity of $\nu^*$ verifies the arrow action. Equations
\eqref{eq:base} make the single arrow $\theta$ the component at every
$t$ of a natural transformation from the constant functor at $(U,a,b)$
to $\mathcal P$. Naturality is $1_S\theta=\theta$.

\begin{corollary}\label{cor:classical}
Under the preceding hypotheses, every pair $(t,\nu^*t)$ with
$t\in\mathcal E(t_0)$ is inseparable. In particular, no such
extension is complete.
\end{corollary}
\begin{proof}
Apply Proposition~\ref{prop:transport} to the component $\theta$.
A complete consistent pair would be its own complementary separator.
\end{proof}

The data in \eqref{eq:base} are already supplied by the one-node
proofs in \cite[Corollary 6.3]{General}. Explicitly, let
$m_A:A\to U$ and $m_B:B\to U$ split the domains $a,b$,
let $c:D\to S$ be the total proof-conclusion map, and let
$t_0=\ran(c)$. The total one-node proof maps satisfy
\[
 \begin{tikzcd}[column sep=large]
 A\arrow[r,"u_A"]\arrow[d,"m_A"']&D\arrow[d,"c"]&
 B\arrow[l,"u_B"']\arrow[d,"m_B"]\\
 U\arrow[r,"\theta"']&S&U\arrow[l,"\nu\theta"]
 \end{tikzcd}
\]
Since $t_0c=c$, these equations give
$t_0\theta m_A=\theta m_A$ and
$t_0\nu\theta m_B=\nu\theta m_B$.
Consequently $m_A^*\theta^*t_0=1_A$ and
$m_B^*\theta^*\nu^*t_0=1_B$, which are exactly
\eqref{eq:base} under the splittings. Thus every extension uses the
same bounded-lattice homomorphism $\theta^*:\D(S)\to\D(U)$.

For an unsplit assertion domain $s\in\D(V)$, take
$\nu:V\to V$ with $\bar\nu=s$ and $s\nu=\nu$, and put
$\nu^*t=\overline{t\nu}$ for $t\le s$.
The corresponding source lattice is the interval
$\D(V)\mathord\downarrow s$, with top $s$.
A total $\theta:U\to V$ satisfying $s\theta=\theta$
has $\theta^*s=1_U$. Splitting $s$ gives the formulation above;
in the unsplit formulation completeness means $t\vee\nu^*t=s$.
The proof in \cite[Theorem 6.1]{Domains} uses this relative form.

The construction of classical syntax and a consistent base theorem
domain remains the substantive premise of this application. The
finite-stage constructions and consistency proofs in
\cite{General,Iteration} establish it under their stated hypotheses;
the three-point computational cosheaf category in \cite{ThreePoint}
is a nonnumerical instance satisfying these hypotheses.
Corollary~\ref{cor:classical}
strengthens the conclusion to inseparability of every extension's
theorem/refutation pair without adding an interpretation hypothesis.

\section{Functors between computational categories}

The preceding reductions are total arrows inside a fixed category.
For a covariant restriction functor $F:\mathscr C\to\mathscr B$,
the induced domain map is covariant:
$F_X:\D_{\mathscr C}(X)\to\D_{\mathscr B}(FX)$,
$e\mapsto Fe$.

\begin{proposition}\label{prop:reflection}
Suppose the domain posets of $\mathscr C$ and $\mathscr B$ have
binary joins and $F$ preserves restriction zeros and these joins.
For a disjoint source pair $(a,b)$, inseparability of $(Fa,Fb)$
implies inseparability of $(a,b)$.
If, in addition, the map
$\mathscr C(X,X)\to\mathscr B(FX,FX)$ is bijective, then $F$
also preserves inseparability of pairs on $X$.
\end{proposition}
\begin{proof}
The functor preserves meets, both bounds, and the domain order.
It therefore sends each source complementary separator to a target
complementary separator, proving reflection.

For preservation, lift a target complementary separator $(r,s)$ to
endomorphisms $h,k$ of $X$, and put $e=\bar h$, $d=\bar k$.
Then $Fe=r$ and $Fd=s$. Injectivity on endomorphisms reflects the
equations $ed=0$, $e\vee d=1$, $ea=a$, and $db=b$ from their
images. Hence $(e,d)$ is a source separator.
\end{proof}

We use Heller's 1990 definition: a recursion category is a
prodominical isotype with a Turing morphism
\cite[\S9, pp.~1264--1265]{Heller}.

\begin{proposition}\label{prop:projection}
Preservation of a universal evaluator, products, coproducts, ranges,
and iteration does not by itself imply preservation of a specified
inseparable pair by a functor between recursion categories.
\end{proposition}
\begin{proof}
Take the numerical recursion category $\mathscr C$ on formal objects
generated from $\mathbb N$ by binary products, binary coproducts, and
nonempty finite words, with partial recursive arrows on their usual
effective presentations. Fix inseparable domains $a,b$ on $\mathbb N$.
Let $\mathscr C^{(2)}$ be the full subcategory of
$\mathscr C\times\mathscr C$ on objects $(X,X)$.
Thus its arrows are pairs $(f,g)$ of arrows of $\mathscr C$ with
the same source and target. Equality, identities, composition,
restrictions, zeros, and ranges are componentwise. Products,
coproducts, and word objects are the corresponding diagonal objects,
with their componentwise structural maps. The numerical range of $f$
is the partial identity on its c.e. image. Agreement after $f$ implies
agreement on this image; hence the full range agreement law, as well
as its product compatibility, holds componentwise.
Finite unions are componentwise. A finite-run iteration in
$\mathscr C^{(2)}$ is the pair of the numerical finite-run iterations:
each component terminates exactly on its own finite-termination domain
and has the corresponding terminal value. Thus iteration and its
termination domains are preserved by either projection. This uses no
arbitrary countable joins.

All objects are isomorphic to $(\mathbb N,\mathbb N)$ by pairs of
the effective coding isomorphisms. If
$\tau:\mathbb N\times\mathbb N\to\mathbb N$ is the numerical
universal evaluator, then $(\tau,\tau)$ is universal in
$\mathscr C^{(2)}$: a pair of computations has a pair of total
program-parameter maps, obtained separately from $\tau$.
The partial-product equations and the zero laws hold componentwise.
Together with isotypy and this evaluator, they make
$\mathscr C^{(2)}$ a recursion category.

Projection to the second component is a full functor
$F:\mathscr C^{(2)}\to\mathscr C$ preserving all the displayed
structures. Let $p,q$ be the domains of the even and odd natural
numbers. The domains $(a,p)$ and $(b,q)$ are disjoint and inseparable in
$\mathscr C^{(2)}$, since the first components of a separator
would separate $a,b$ in $\mathscr C$.
Their images are the nonzero complementary pair $p,q$.
\end{proof}

This source is not dominical: the nonzero arrows $(1,0)$ and $(0,1)$
have zero restriction product, violating the additional dominical
condition of \cite[Proposition 5.1(ii)]{Heller}.
The projection is not faithful and does not reflect zero. Its target
still has the inseparable pair $(a,b)$; the counterexample concerns
the image of the specified source pair. Reflection in
Proposition~\ref{prop:reflection} likewise concerns a pair with specified
source preimages. Neither statement identifies arbitrary target pairs
with source pairs.

\begin{thebibliography}{9}
\bibitem{Domains} F. Lengyel, \emph{Domains, finite proofs, and essential
incompleteness}, 29 October 2026, 23-page manuscript.
\bibitem{General} F. Lengyel, \emph{Finite proof algebras in range
restriction categories}, 29 October 2026, 13-page manuscript.
\bibitem{Iteration} F. Lengyel, \emph{Iteration domains, finite proofs,
and computational cosheaves}, 29 October 2026, 14-page manuscript.
\bibitem{ThreePoint} F. Lengyel, \emph{Range completion and computational
cosheaves on a three-point space}, 29 October 2026, 7-page manuscript.
\bibitem{Heller} A. Heller, An existence theorem for recursion categories,
\emph{Journal of Symbolic Logic} \textbf{55}(3) (1990), 1252--1268.
\href{https://doi.org/10.2307/2274486}{doi:10.2307/2274486}.
\end{thebibliography}
\end{document}
